\scanpage{27}
% L01–06
Les réductions habituelles nous ramènent au cas d'un anneau \inserted{de base} \deleted{artinien $\mathbb Z/p^n\mathbb Z$ $A$, corps} artinien $A$, corps résiduel \deleted{\unsure{$k=$|$k$}} $\mathbb Z/p\mathbb Z$, à prouver que si $\lambda\deleted{\unsure{\ne|=}}{,}\lambda'\in A^*$ sont tels que $\lambda+\lambda^{-1}=\lambda'+\lambda'^{-1}$, alors on a $\lambda'=\lambda^{\deleted{\unsure{\varepsilon}}}$ ou $\lambda'=\lambda^{-1}$. Supposons que l'on sache $\lambda'_n=\lambda_n$ \unsure{prouvons|prouver}
\[
\underset{\substack{\Vert\\\lambda'}}{\lambda'_{n+1}}=\underset{\substack{\Vert\\\lambda}}{\lambda_{n+1}}^{\pm1}\ ?
\]

% L07–11
On a \deleted{$\unsure{\lambda'_n=\lambda_n}$} $\lambda'=\lambda+\mu$, avec $\unsure{\mathfrak m|m}\mu=0$, \deleted{d'où} donc
\[
\begin{aligned}
\lambda'&=\lambda(1+\lambda^{-1}\mu)=\lambda(1-\mu)=\lambda+\mu\\
\lambda'^{-1}&=\lambda^{-1}\underbrace{(1-\mu)^{-1}}_{1+\mu}=\lambda^{-1}(1+\mu)=\lambda^{-1}\unsure{-|+}\mu\\
\lambda'+\lambda'^{-1}&=\lambda+\lambda^{-1}+\mu(\lambda^{-1}-\lambda^{\deleted{\unsure{-1}}})=\lambda+\lambda^{-1}.
\end{aligned}
\]
% slanted
\margin{\uline{NB} Il s'agit apparemment d'une propriété générale des morphismes \deleted{\unsure{qui}} ``ramifiés'' (non nets) en un point : \unsure{certaine|certains|existence} ``\unsure{continu|continuum}'' de \unsure{l'ens.} des solutions infinitésimales (dès le $1^{\rm er}$ ordre\ldots)}

% L12–16
Donc l'hypo $\lambda'+\lambda'^{-1}=\lambda^{\deleted{\unsure{-1}}}+\lambda^{-1}$ résulte de $\lambda'_n=\lambda_n$ i.e. $\lambda'=\lambda+\mu$, $\mathfrak m\mu=0$ \unsure{$\}$|?} ; a-t-on pour autant
\[
\lambda'=\lambda\quad\text{ou}\quad\lambda'^{-1}=\lambda\quad\deleted{\unsure{?}}?
\]
\deleted{i.e} \unsure{Or|On} ces relations signifient, \deleted{soit} \inserted{l'une} $\mu=0$,

% L17–23
\deleted{soit} \inserted{l'autre} \deleted{$\unsure{\lambda'=\lambda(1+\mu)=\lambda|\lambda^{-1}(1+\mu)=\lambda}\quad\text{i.e.}\quad\lambda\mu=0$}
\[
\begin{aligned}
&\deleted{\lambda^{-1}(1+\mu)}\quad\deleted{\lambda-\lambda^{-1}+\mu=0}\\
&\mu=\lambda^{-1}-\lambda
\end{aligned}
\]
donc ne correspondent qu'à deux valeurs possibles \inserted{au plus} de $\mu$ -- donc (si $k$ n'est pas réduit à $2$ él) on peut choisir $\mu$ de façon à mettre en défaut la relation $[(\lambda'=\lambda)\text{ ou }(\lambda'=\lambda^{-1})]$.
% archive sheet number: 26
